\section{Discussion}

The central difficulty of this study is that sequencing depth and tissue density are, in
this cohort, the same measurement. A domain's marker stability, the apparent
stage-dependence of the spatial programmes, the cohort-level copy-number difference and the
prognostic weight of one domain's signature all vary with depth, and depth varies with the
amount of tumour in a spot. Adjusting for depth is consequently not a neutral correction: it
removes part of the biological effect together with the technical one. Adjusted estimates
are conservative and may be over-corrected; unadjusted estimates are inflated. Where the
choice is consequential, both are reported rather than one being nominated as correct.

The coupling also accounts for the behaviour of the pathology model, which is the one
result here obtained without any transcriptomic measurement. Its apparent discrimination
was largely a measure of how much tissue a spot contained, and on the axis this study needs,
the separation of precursor from invasive lesions, it failed outright. A model that reads
tissue density will appear to succeed on this disease for as long as density and malignancy
coincide, which is exactly the condition that makes the confound difficult to detect.

This also accounts for the narrower scope of some results relative to the source study.
Where a claim depends on comparing dense with sparse tissue, which in this disease means
comparing invasive with precursor tissue, the design carries the confound in its structure.
The only clean comparisons are those that hold depth fixed, either by pairing within a
section or by matching across sections. That is why the one copy-number result retained is a
within-section comparison, and why the disease-axis queries were constructed by pairing
within patients.

\subsection{What the positive results establish}

Four findings survive the checks applied. First, the tissue resolves into seven recurring
niche domains, six of which are ordinary lung structures and only one of which is confined
to invasive disease. The lymphoid-aggregate domain is present in precursor stages, which
argues against attributing immune aggregation to the onset of malignancy. Second, the
ECM/interstitial programme that distinguishes the invasive domain is itself elevated in
invasive disease within patients (20 of 22), and the result holds when the genes added by
hand are removed and under every leave-one-out variant, so the domain label does not rest on
a hand-picked panel. Third, the spatial copy-number signal lies in the AT2 compartment and
is present before invasion, consistent with AT2 as the cell of origin and with an orthogonal
expression-based classifier. Fourth, a query signature derived from a patient-internal
disease axis and filtered on independent survival evidence returns a reproducible set of 34
compounds, dominated by PI3K/mTOR inhibitors, which act on the ECM/interstitial and AT2
domains rather than on the tissue as a whole.

The candidate list should be read as a hypothesis about a compartment rather than as a
therapy. Several caveats are structural. The compounds act in cancer cell lines and have not
been tested in lung fibroblasts or in any primary tissue. The disease axis against which they
were selected is a tissue-level average, so the candidates reverse the aggregate state and
need not act on any single cell type. The 34 compounds are the intersection of three scoring
schemes, but those schemes share a query and a library and their agreement is therefore a
robustness check rather than three independent measurements. The split-half reproducibility
of 0.962 describes the stability of the ranking under resampling of patients and says
nothing about whether the ranking is biologically correct. No compound reported here has
been tested experimentally.

\subsection{Methodological observations}

Several of the checks applied here are reusable, and the difficulties encountered carry
lessons that are not specific to this dataset.

A positive control with a known answer, run through the identical pipeline, distinguishes
``the method found nothing'' from ``the data contain nothing detectable by this method''.
Without such a control the negative results reported above would be uninterpretable. A
related practice is to inspect the distribution of a summary statistic rather than the
summary alone: an AUROC of exactly 0.5000 in this study was a degeneracy, all cells carrying
one identical value, and not a true null.

Control for a nuisance variable is better served by sensitivity analysis than by a single
adjusted estimate. When a measurement is strongly coupled to such a variable, examining
subsets of the measurement is more informative than adjusting once. In this study, showing
that a result held when the eight hand-added genes of a marker panel were removed, and
across all twenty-two leave-one-out variants, established that the finding did not depend on
the panel's arbitrary members. A related caution applies to sub-clustering: because clusters
are selected on the same data used to test them, per-subtype contrasts carry a
selective-inference penalty that is not corrected here \citep{gao2024selective}, which is
one reason the per-compartment gene list is reported as descriptive.

Two further observations concern how criteria are constructed and reported. A gate that
every configuration passes is not a gate. The spatial-coherence criterion applied here,
which required a partition's coherence to exceed a label-permutation null on at least 90\%
of sections, was passed by essentially every configuration tried, including one that uses no
spatial information at all ($\lambda=0$; coherence 0.54--0.60 against a null 95th percentile
of 0.13--0.18). Tissue expression carries spatial autocorrelation intrinsically, so a purely
expression-based partition clears the null by a wide margin. The criterion can exclude a
grossly anti-spatial partition but cannot demonstrate that spatial information was used;
ranking was accordingly decided by cross-seed stability instead. In addition, a criterion
introduced by the analysts must not be presented as though it carried external authority.
The coherence gate is ours. The source study sets no spatial gate, and the BANKSY authors'
recommendations \citep{singhal2024banksy} concern parameters ($\lambda$, $k_{\text{geom}}$,
resolution) rather than validation criteria.

Finally, a number produced by an optimiser's trajectory is not a measurement. One weighting
parameter carried for a time in this work, $w_{\min}=(1/6)\times0.7^{k}$, is the path the
solver took rather than a quantity estimated from the data, and should not be reported as
though a fit had produced it.

\subsection{Relation to the source study}

This study set out to extend a published description of the disease sequence
\citep{peng2026}. The results agree with it on the identity of the niche structures and on
AT2 as the compartment carrying copy-number signal, and differ in emphasis on three points.
The source study reports a pro-inflammatory niche that is enriched early and less frequent
in invasive disease; our spatial arm does not reproduce that direction. Its developmental
trajectory places the least differentiated state in precursor lesions, whereas ours runs the
other way, although, as noted, the measurement is ambiguous and is not treated as a
refutation. The source study also does not describe ambient-RNA correction, whereas one of
its samples is, in our hands, dominated by ambient signal. These divergences are reported
rather than adjudicated, since each is confounded by a design difference: different
platforms, a stage axis rather than a subtype axis, and different approaches to quality
control.

\subsection{Priorities for further work}

The comparison that would remove the dominant confound is within a single section, comparing
the invasive front with regions of the same section distant from it and holding depth fixed
by construction. That analysis was not performed and is the one we would prioritise next,
together with testing the candidate compounds in primary lung fibroblasts or in organoids
derived from the same patients.
